\section{模型的评价与推广}

\subsection{模型的优点}

\textbf{第一，单一内核覆盖四个问题，口径可追溯。} 四个问题共用同一套 10 分钟粒度的线性规划内核，差异被完全参数化到时间、信息、结算与价格四个维度上（表~\ref{tab:param}）。变量定义、单位、参数值与时间范围在四问之间逐项一致，因此任何一个问题的实现都可以用同一套核验脚本检查，也使得“净负荷”“弃光量”这类与价格无关的量成为跨问题的硬校验点：四问的交付期净负荷均为 17\,939\,189.494967 kWh，问题二与链 \emph{4-2} 的弃光量均为 990\,168.090312 kWh，而费用必然不同——这既证明价格情景真实生效，也排除了“把上一问结果填到下一问”的可能。

\textbf{第二，实现正确性被当作一等目标，并且给出了硬证据。} 本文构造了三类相互独立的验证手段：一是\textbf{解析界}，为每个问题独立构造可行上界与下界（如问题一的“不用储能”策略给出 48052.046591 元），并明确指出因附件 4 存在 0.0076 元/kWh 的极端低值、以最小电价构造的下界不具判别力；二是\textbf{独立第二实现}，用不同的变量排序、变量消去与求解器设置重算，问题二中与主模型逐位相同；三是\textbf{退化情景锚点}，把链 \emph{4-3} 的价格换成附件 1 的单日曲线后必须逐位复现问题三的六个数值，实测相对差不超过 $4.86\times10^{-13}$（判据为 $10^{-9}$）。这三类手段把“实现是否正确”与“模型是否合适”两个问题彻底分开。

\textbf{第三，量纲与索引口径被显式固定，避免了最隐蔽的一类错误。} 本文明确规定：决策变量以 kWh 为单位、本身已含时间因子，费用表达式中不再乘 $\Delta t$；附件时间戳为所属时段的左端点，据此计划窗为 0:10 至 24:10。这两项约定看似细节，但一旦出错，费用会整体偏差 6 倍、全部时段整体错位，而结果“看起来仍然合理”。核验中对每一项都做了独立复算：若误乘 $\Delta t$，问题二交付期费用会落到约 $2.04\times10^{6}$ 元，与正确值相差 6 倍。

\textbf{第四，对“未通过”与“不利结果”的处理方式本身构成结论的一部分。} 四个问题共登记六项未通过判据与反例（表~\ref{tab:rob} 及相关正文），全部如实保留并逐条给出归因与适用边界。例如问题二的“脆弱”分级并非模型失败，而是准确指出了日尺度量对效率口径的高敏感性；链 \emph{4-3} 的“条件稳定”则准确指出了模型对光伏侧误差的暴露远大于价格侧——这一结论对实际运行有直接指导意义。

\textbf{第五，计算预算可控且不依赖特殊硬件。} 最大规模实例为问题二的 52\,560 个时段、315\,360 个变量、105\,120 条等式约束，在普通 CPU 上以稀疏矩阵求解约 5.8 s；问题三与链 \emph{4-3} 因分解为 1460 个单日层分别约需 42 s 与 39 s；全部鲁棒性与消融实验（合计上千次线性规划）在数十分钟内完成。这意味着模型可以被反复重算，适合作为决策支持工具\cite{ref-hai-2025-microgrid-price-ems}。

\subsection{模型的不足与适用边界}

\textbf{第一，确定性口径，预报误差的代价只能事后量化。} 主结果不含误差分布，也不含随机或鲁棒优化。问题三比问题二贵 2\,307\,565.50 元（$+18.863\%$），同净负荷下这一差额完全由预报误差造成；链 \emph{4-3} 比链 \emph{4-2} 贵 2\,282\,639.673590 元。本文用“价格预报误差的代价”（520\,790.839035 元与 258\,485.465704 元）与“光伏预报误差的代价”（紧急购电费 1\,659\,896.438588 元）分别量化了这两类误差，但都属事后度量；消融中的随机对照说明改进空间在 3\% 量级，不足以为更换主口径提供充分理由。

\textbf{第二，光伏降尺度不保块内能量。} 附件 3 只给整点预报，本文按整点锚定线性插值降到 10 分钟，块内平均功率为 $(2.5A_{k-1}+3.5A_k)/6$，一般不等于整点值。若改用分段常数前向保持，费用会上升 15.83\%，说明该项口径影响很大。本文的做法是保留双端整点信息并显式披露其能量性质，但这是本问题最需审慎的一项建模选择。

\textbf{第三，若干物理简化与题面参数的解释约定。} 模型假设常效率（不随荷电状态、功率与温度变化）、无自放电、购电无功率上限、无售电收益、不计储能折旧；其中“5000 kW 同时作用于并网点侧与电池侧、取更严者”是本文对题面参数的解释约定，不构成文献支持的结论。这些简化的影响方向已被量化：效率是唯一一阶敏感参数；若存在 3000 kW 的隐含购电上限，问题一的计划将不可行；在允许售电的情景下，松散的弃光上界会出现售出电量超过物理盈余的非物理套利。

\textbf{第四，文献证据等级不足。} 全文引用的文献均只达到摘要或元数据级，没有一条经过全文精读，因此引用只用于支撑框架级与机制级主张（“这类问题可以用线性规划建模”“分时电价下存在充放电套利动机”“偏差结算制度会影响储能策略”），不用于支撑任何公式级或定量级结论。题面给定的硬参数（紧急购电 5 倍价、偏差结算 0.5 倍与 1.5 倍）与本文的模型内推论（紧急购电恒零）同样不作为文献结论陈述。

\textbf{第五，最优面非唯一带来呈现上的技术限制。} 线性规划的最优调度面可能不唯一（储电量上下限两端均可能活跃），因此验收以约束残差、结构恒等式、目标值与交付表格为准，不比对逐点解的唯一性；本文统一采用字典序规则选取提交解，并量化了该规则的影响（问题三约 $10^{-4}$ 量级、链 \emph{4-3} 为 0.110739 元），说明它不改变任何结论，但读者不应把不同求解器返回的轨迹差异理解为模型不稳定。

\textbf{第六，前瞻深度与日边界行为的结构依赖。} 在“逐日 0:00 决策加终端自由”的组织方式下，四条链的日末储电量都贴住下限 1200 kWh。滚动时域对照证明这只是该结构的产物：把前瞻深度提高到 3 天以上，日末储电量均值即升至约 5330 kWh。因此本文不把“日边界不携带电量”写成普遍规律——这既是结论，也是对模型适用边界的提示。

\subsection{模型的改进方向}

其一，\textbf{提升证据等级}：取得关键文献全文，把引用从摘要级提升到全文级，重新核验公式级与定量级主张，尤其是偏差结算制度与价格预测方法相关的主张。

其二，\textbf{补充误差分布的实测证据}，在此基线上重新评估随机规划与分布鲁棒优化是否应作为主口径。本文已给出评估框架（两阶段随机、滚动时域、二者组合）与预注册判据，但结论依赖场景构造，需要价格与光伏误差的真实分布数据支撑。

其三，\textbf{细化储能物理模型}：引入与荷电状态、充放电功率相关的动态效率与自放电项，并在此基线上重估费用与稳定性分级；本文已量化效率的敏感性，可作为该改进价值的前置估计。

其四，\textbf{引入极短的实时平衡层}：题面允许在“其他时刻”获得预报，本文结论限于现有发布时刻下的决策频率；若能获得更频繁的预报发布，则需重新评估“信息供给增加”的价值。

\subsection{本文的创新点}

\textbf{创新点一：以“四个维度上的参数化”统一四个问题，并由此得到一个可复用的建模模板。} 本文没有为四个问题分别建模，而是先证明它们共享同一物理系统与同一数学结构，再把差异定位到时间、信息、结算、价格四个维度。这既减少了重复推导，更让“什么是这道题的难点”变得可判定：问题一、二在时域与终端口径，问题三在信息结构与降尺度，问题四在价格可观测性。

\textbf{创新点二：给出互补性引理，把“是否需要二元变量”从工程判断变成数学结论。} 通过对往返效率小于 1 的利用，本文证明同时充放电被严格支配，从而线性松弛无损；该结论被两个不同规模的整数模型对照独立验证（问题一中 144 个二元变量模型目标值完全相同，问题二中 52\,560 个二元变量模型仅差 $3.2\times10^{-7}$ 而耗时增加 50 倍）。这节省的不只是计算量，也避免了因引入大 M 常数而带来的数值脆弱性。

\textbf{创新点三：把“代价”做成可分解的量，而不是笼统的“不确定性影响”。} 本文对两类误差分别给出价格与金额，并在问题四中把两条链 2\,282\,639.673590 元的总差额完全分解为计划口径差、偏差结算与紧急购电三项（恒等式残差不超过 $10^{-6}$ 元）。这种分解使“哪一条通道更值得改进”有了定量答案：对链 \emph{4-3} 而言，降低光伏预报误差的收益远大于提高价格预测精度。

\textbf{创新点四：用“退化情景复现”构造实现正确性锚点。} 把链 \emph{4-3} 的价格换成附件 1 的单日曲线后，模型应当精确退化为问题三；实测六项指标逐位复现，相对差不超过 $4.86\times10^{-13}$。这一手法把“实现是否正确”变成一个可以客观判定的问题，且适用于任何“旧问题加新情景”的题目结构。

\textbf{创新点五：用滚动时域对照回答“结构性问题”，而非仅作灵敏度扫描。} “日边界储电量为什么总贴下限”无法通过参数扰动回答，因为它是结构性的。本文设计了对决策时刻与承诺规则完全不变、只改变“优化时看多远”的对照实验：前瞻深度为 1 天时逐位复现主口径，提高到 3 天以上时日末储电量均值升至约 5330 kWh、贴下限日数由 364/365 降至 30/365。这一“只改一个维度、其余严格不变”的设计使结论具有因果含义。

\textbf{创新点六：对不确定性与前瞻给出“三层递进回答”，把建议落到可执行的决策上。} 本文对“是否需要更多预报”给出层次分明的答案：需要调整机制、不需要加密决策时刻；需要一定前瞻（至少 3 天）、不需要无限前瞻；随机优化有收益（3.29\%）但幅度不足以替代主口径。同时指出优先级：由于模型对光伏侧误差的暴露远大于价格侧，提高光伏预报精度应是首要投入。
